\documentclass[11pt]{article}
\usepackage[T1]{fontenc}
\usepackage{amsmath,amssymb,amsthm}
\usepackage[margin=1in]{geometry}
\usepackage{booktabs}
\usepackage{hyperref}
\hypersetup{
  hidelinks,
  pdftitle={The Game Chromatic Number of the Thirteen-Vertex Graph M3(C3)},
  pdfauthor={Xiaojian Zeng},
  pdfsubject={Exact computer-assisted determination of a finite graph-coloring game}
}

\newtheorem{theorem}{Theorem}
\newtheorem{lemma}[theorem]{Lemma}

\newcommand{\Zthree}{\mathbb Z/3\mathbb Z}
\newcommand{\chig}{\chi_{\mathrm g}}

\title{The Game Chromatic Number of the Thirteen-Vertex Graph
$M_{3}(C_{3})$}
\author{Xiaojian Zeng\\
PBC School of Finance, Tsinghua University, 100190\\
\texttt{xiaojianzeng@foxmail.com}}
\date{}

\begin{document}
\maketitle

\begin{abstract}
We determine the game chromatic number of the fixed generalized-Mycielski
graph $M_{3}(C_{3})$.  An elementary propagation argument excludes three
colors.  For four colors, an exact minimax computation produces a complete
Alice-strategy directed acyclic graph, and a separately implemented checker
reconstructs the graph and its symmetries, replays every stored move, exhausts
every raw Bob reply, checks both terminal conditions, and performs backward
induction.  The certificate has $22{,}071$ canonical positions and proves
that Alice wins.  Consequently
\[
                 \chig\bigl(M_{3}(C_{3})\bigr)=4.
\]
The result is only for the layer parameter $3$ and makes no assertion about
the remaining generalized layers.
\end{abstract}

\section{The fixed graph and the coloring game}

The graph-coloring game was introduced in graph form by
Bodlaender~\cite{Bodlaender1991}.  Alice and Bob alternately color a previously
uncolored vertex from a common palette, with Alice moving first and every
intermediate coloring required to be proper.  Alice wins if the whole graph
is colored.  Bob wins as soon as an uncolored vertex has no legal color.  The
game chromatic number $\chig(G)$ is the least palette size for which Alice has
a winning strategy.

We use the generalized-Mycielski notation of Mou, Sun, and
Zhang~\cite{MouSunZhang2026}.  For the present paper it is enough to define the
target graph directly.  Let $H=M_{3}(C_{3})$ have vertex set
\[
 V(H)=\{v_i^m:i\in\Zthree,\ 0\leq m\leq3\}\cup\{\omega\}.
\]
Its edges are
\begin{align*}
 &v_i^0v_j^0 &&(i\ne j),\\
 &v_i^mv_j^{m+1} &&(0\leq m\leq2,\ i\ne j),\\
 &v_i^3\omega &&(i\in\Zthree).
\end{align*}
Thus $H$ has $13$ vertices and $24$ edges.  The nine vertices in layers
$0,1,2$ have degree $4$, while the three vertices in layer $3$ and $\omega$
have degree $3$.

Theorem~5.1(1) of Mou--Sun--Zhang gives
$4\leq\chig(M_k(C_3))\leq5$ for every $k\geq2$, but does not choose between
the two values for $k=3$~\cite{MouSunZhang2026}.  The contribution here is the
exact value for that one thirteen-vertex instance.  In particular, we do not
claim that the source paper explicitly posed this individual instance as an
open problem.

\begin{theorem}\label{thm:main}
For the graph $H=M_3(C_3)$ defined above,
\[
                         \chig(H)=4.
\]
\end{theorem}

\section{The lower bound}

\begin{lemma}\label{lem:lower}
The graph $H$ is not three-colorable.  Hence $\chig(H)\geq4$.
\end{lemma}

\begin{proof}
Suppose that a proper three-coloring exists.  The base layer is a triangle, so
write $c_i$ for the color of $v_i^0$; the three colors $c_0,c_1,c_2$ are
distinct.  The vertex $v_i^1$ is adjacent to the two base vertices whose
indices differ from $i$, and therefore its only available color is $c_i$.
The same argument propagates: if $v_i^{m-1}$ has color $c_i$ for every $i$,
then $v_i^m$ is adjacent to the other two colors and must have color $c_i$.
Consequently the three neighbors $v_0^3,v_1^3,v_2^3$ of $\omega$ use all
three colors, leaving no color for $\omega$.  This contradiction proves
$\chi(H)\geq4$, and $\chi(H)\leq\chig(H)$ gives the assertion.
\end{proof}

\section{A proof-carrying four-color strategy}

We now describe the exact finite certificate for the upper bound.  Order the
vertices as
\[
 v_0^0,v_1^0,v_2^0,v_0^1,\ldots,v_2^3,\omega
\]
and encode a position by a length-$13$ word over
$\{0,1,2,3,4\}$, where $4$ denotes an uncolored vertex.  All comparisons in
the search and replay are integer and exact.

Every permutation of the three indices, applied simultaneously in all four
layers and fixing $\omega$, is an automorphism.  These are all the
automorphisms of $H$.  Indeed, $\omega$ is the unique degree-$3$ vertex having
three degree-$3$ neighbors, so it is fixed.  Distance from $\omega$ then fixes
each layer set.  Between two consecutive layers, the unique missing partner
of $v_i^m$ is $v_i^{m+1}$; hence a permutation of the base indices forces the
same permutation in each later layer.  The automorphism group therefore has
order $6$.

The search identifies positions under these six automorphisms and under the
$24$ global permutations of the four color names.  Its canonical
representative is the lexicographically least word in the resulting orbit.
The certificate is a directed acyclic graph of canonical positions.  At an
Alice position it stores one legal move.  At a Bob position it stores every
raw legal vertex--color pair, even when several such pairs have the same
canonical child.  A stored terminal position is accepted only if all thirteen
vertices are properly colored; a position with an uncolored vertex seeing all
four colors is instead a Bob terminal.

The following elementary observation is the logical core of the
computer-assisted proof.

\begin{lemma}[certificate soundness]\label{lem:soundness}
Suppose a finite directed acyclic graph rooted at the empty position has the
properties just stated: every terminal is a fully colored Alice win; every
nonterminal Alice node has a stored legal edge; and every nonterminal Bob node
contains every raw legal reply.  If every edge leads to the canonical orbit
representative of the position obtained by its move, then the root is winning
for Alice.
\end{lemma}

\begin{proof}
Induct downward on the number of uncolored vertices.  A terminal node is an
Alice win.  At an Alice node she takes its stored edge, while at a Bob node
every move available to Bob occurs among the stored edges.  Thus, if all
children are already certified for Alice, so is the parent; the induction
reaches the root.

Canonicalization does not change this conclusion.  A graph automorphism and
a global color permutation preserve properness, legal moves, whose turn it
is, and both terminal conditions.  At an actual position Alice transports the
stored move from its canonical representative back through the same symmetry.
After either player's move, the resulting position can again be transported
to the stored canonical child.  Hence the backward induction on orbit
representatives defines a strategy in the original game.
\end{proof}

Separating discovery code from a reconstruction-based checker is a standard
proof-carrying pattern for large exact finite computations; for a recent
example using exhaustive symmetry reduction and an independently replayed
certificate, see Zeng, Liu, Ratnavelu, and Ong~\cite{ZengEtAl2026Certificates}.
Here the generator uses first-occurrence color normalization after applying
each graph automorphism.  The checker does not import the generator.  It
reconstructs the graph from the displayed edge formula, enumerates graph
automorphisms using NetworkX's isomorphism machinery~\cite{HagbergSchultSwart2008},
and obtains exactly six.  Independently of the generator's normalization, it
then tries all $6\cdot24$ graph/color symmetries to canonicalize every child.

The checker verifies the following conditions throughout the certificate:
\begin{itemize}
  \item every stored position is proper and canonical;
  \item every edge is a legal move, colors exactly one new vertex, and names
        the correct canonical child;
  \item every Bob node contains each raw legal vertex--color pair exactly
        once, without orbit deduplication;
  \item terminal labels agree with the two game-ending conditions; and
  \item bottom-up evaluation certifies Alice at the empty root.
\end{itemize}
The resulting counts are recorded in Table~\ref{tab:certificate}.  As a small
boundary check, the prescribed first move colors $v_0^0$ with color $0$; the
next node contains all $44$ legal raw replies available to Bob.

\begin{table}[ht]
\centering
\begin{tabular}{lr}
\toprule
quantity & value\\
\midrule
canonical positions in the strategy DAG & $22{,}071$\\
directed edges & $61{,}173$\\
raw Bob-response edges & $45{,}761$\\
terminal positions & $400$\\
graph automorphisms independently enumerated & $6$\\
color permutations tested per automorphism & $24$\\
\bottomrule
\end{tabular}
\caption{Exact certificate and replay counts.  All $400$ stored terminals are
fully colored Alice wins.}
\label{tab:certificate}
\end{table}

The successful replay establishes every hypothesis of
Lemma~\ref{lem:soundness}; therefore Alice wins with four colors and
$\chig(H)\leq4$.  Combining this with Lemma~\ref{lem:lower} proves
Theorem~\ref{thm:main}.

\section{Reproducibility, scope, and assistance}

The supplementary files are \texttt{solve\_game.py},
\texttt{q4\_strategy.json.gz}, \texttt{verify\_q4\_strategy.py}, and the
machine-readable report \texttt{q4\_verification.json}.  From the project
root the certificate and replay report are regenerated by
\begin{verbatim}
python3 evidence/solve_game.py --output evidence/q4_strategy.json.gz
python3 evidence/verify_q4_strategy.py evidence/q4_strategy.json.gz \
  --report evidence/q4_verification.json
\end{verbatim}
The generator is deterministic and the verifier performs no floating-point
comparisons.  The proof determines only $M_3(C_3)$.  It does not generalize
the four-color strategy to $M_k(C_3)$ for any other $k$, nor does it establish
a priority claim beyond the specific literature comparison made above.

Python programs were used for exact minimax search, certificate generation,
graph-isomorphism enumeration, and exhaustive replay.  Generative-AI systems
assisted with problem triage, program development, audit-oriented proof
organization, literature searching, and manuscript drafting.  No theorem is
accepted on the basis of model output: the finite upper bound is carried by
the explicit certificate and exact replay checks, and the author is
responsible for the statement, scope, and presentation.

\bibliographystyle{plain}
\bibliography{references}

\end{document}
